\section{Fondamenti teorici}
\label{sec:theory}

Si riportano solo i concetti necessari a comprendere le implementazioni. La
notazione segue Sutton \& Barto~\cite{sutton2018}.

\subsection{MDP, valore e policy}
Un problema di RL \`e modellato come un \emph{Markov Decision Process} (MDP):
ad ogni \emph{timestep} l'agente osserva uno \emph{stato} $s$, sceglie
un'\emph{azione} $a$ secondo una \emph{policy} $\pi(a|s)$, riceve un
\emph{reward} $r$ e transita in $s'$. L'obiettivo \`e massimizzare il
\emph{return} atteso $G_t = \sum_{k=0}^{\infty}\gamma^k r_{t+k+1}$, dove
$\gamma \in [0,1]$ \`e il fattore di sconto. Si definiscono la
\emph{state-value} $V^\pi(s)=\mathbb{E}_\pi[G_t\,|\,S_t=s]$ e la
\emph{action-value} $Q^\pi(s,a)=\mathbb{E}_\pi[G_t\,|\,S_t=s, A_t=a]$.

\subsection{Predizione vs controllo}
\`E fondamentale distinguere i due compiti, perch\'e la consegna li richiede
entrambi:
\begin{itemize}[nosep]
  \item \textbf{Predizione}: stimare $V^\pi$ per una policy $\pi$
        \emph{fissata}. Monte Carlo e TD(0) sono algoritmi di predizione.
  \item \textbf{Controllo}: trovare una policy migliore. Q\hbox{-}Learning e
        SARSA sono algoritmi di controllo.
\end{itemize}

\subsection{Monte Carlo (MC) e TD(0)}
Entrambi stimano $V^\pi$ ma con target diversi (S\&B \S5.1, \S6.1):
\begin{align}
  \text{MC:}\quad   & V(S_t) \leftarrow V(S_t) + \alpha\,[\,G_t - V(S_t)\,] \\
  \text{TD(0):}\quad& V(S_t) \leftarrow V(S_t) + \alpha\,[\,R_{t+1} + \gamma V(S_{t+1}) - V(S_t)\,]
\end{align}
MC usa il return \emph{effettivo} $G_t$: \`e non distorto (\emph{unbiased}) ma
ad alta varianza e richiede episodi completi. TD(0) usa il \emph{bootstrap}
$R_{t+1}+\gamma V(S_{t+1})$: \`e distorto (usa una stima) ma a varianza pi\`u
bassa e aggiornabile online. Il termine
$\delta_t = R_{t+1}+\gamma V(S_{t+1})-V(S_t)$ \`e il \emph{TD error}.

\subsection{SARSA (on-policy) e Q-Learning (off-policy)}
Sono aggiornamenti TD sulla $Q$ (S\&B \S6.4--6.5):
\begin{align}
  \text{SARSA:}\quad     & Q(S,A) \leftarrow Q(S,A) + \alpha\,[\,R + \gamma Q(S',A') - Q(S,A)\,] \\
  \text{Q\hbox{-}Learning:}\quad & Q(S,A) \leftarrow Q(S,A) + \alpha\,[\,R + \gamma \textstyle\max_{a'} Q(S',a') - Q(S,A)\,]
\end{align}
SARSA \`e \emph{on-policy}: il target usa l'azione $A'$ effettivamente scelta
dalla policy $\varepsilon$-greedy, quindi tiene conto del rischio
dell'esplorazione. Q\hbox{-}Learning \`e \emph{off-policy}: il target usa il
$\max$, cio\`e valuta la policy greedy indipendentemente da come si esplora.

\subsection{Esplorazione: $\varepsilon$-greedy}
La policy $\varepsilon$-greedy sceglie un'azione casuale con probabilit\`a
$\varepsilon$ (\emph{exploration}) e l'azione greedy $\arg\max_a Q(s,a)$
altrimenti (\emph{exploitation}). \`E il compromesso pi\`u semplice tra i due
e viene usato in tutto il progetto (S\&B \S2.1).

\subsection{Discretizzazione (state aggregation)}
I metodi tabellari richiedono uno spazio di stati finito. CartPole e Acrobot
hanno stato continuo, quindi ogni dimensione viene divisa in \emph{bin}
uniformi (S\&B \S9.4, \emph{state aggregation}): stati vicini finiscono nello
stesso bin. Il numero di stati cresce come (bin)$^{\text{dim}}$: $10^4$ per
CartPole (gestibile), fino a $10^6$ per Acrobot con 10 bin (\emph{curse of
dimensionality}).

\subsection{Deep Q-Network (DQN)}
Quando lo spazio degli stati \`e troppo grande (pixel di Pong), la $Q$ viene
approssimata da una rete neurale $Q(s,a;\theta)$~\cite{mnih2015}. Due
componenti stabilizzano l'addestramento:
\begin{itemize}[nosep]
  \item \textbf{Experience Replay}: le transizioni sono salvate in un buffer e
        campionate a mini-batch casuali, rompendo la correlazione temporale.
  \item \textbf{Target Network}: una copia $\theta^-$ della rete, aggiornata
        di rado, fornisce target stabili
        $y = R + \gamma \max_{a'} Q(S',a';\theta^-)$.
\end{itemize}
Il \textbf{Double DQN}~\cite{vanhasselt2016} disaccoppia selezione e
valutazione dell'azione,
$y = R + \gamma\, Q(S', \arg\max_{a'}Q(S',a';\theta);\theta^-)$, riducendo la
sovrastima ($\emph{maximization bias}$, S\&B \S6.7).
